\section*{Abstract}

The aim of this thesis was to conduct a comparative analysis of monolithic
and microservice architectures in terms of performance, scalability,
and maintenance costs in cloud environments. As software systems become
increasingly complex and the number of users continues to grow, architectural
models that provide effective resource management, flexibility, and reliability
gain significant importance. This work investigates how the choice of software
architecture influences real application performance parameters and what
operational and financial consequences it entails for an organization.

A systematic literature review was carried out to identify current approaches
to measuring the performance and scalability of distributed systems, as well
as existing benchmarking tools. However, most available studies examine these
architectures without a practical cloud deployment context, exposing a research
gap — the lack of a comprehensive comparison of monolithic and microservice
applications deployed in equivalent environments, using unified load-testing
methods and identical test scenarios.

As part of the thesis, a benchmarking application was implemented in two
architectural variants — monolithic and microservice — and both versions were
deployed on a cloud platform. Load tests were conducted using a dedicated tool,
and vertical and horizontal scaling were analyzed. Additionally, the impact of
infrastructure configuration, communication mechanisms, multithreading, and
resource management strategies on overall application efficiency was examined.

The study results indicate that monolithic architecture provides higher
performance under low load and lower maintenance costs in simpler deployments.
In contrast, microservice architecture offers significantly better horizontal
scalability, component independence, and greater flexibility in optimizing
specific parts of the system. At the same time, microservices introduce
additional costs related to inter-service communication, infrastructure
complexity, and operational overhead.

The thesis concludes with a discussion of study limitations, the scientific
contribution of the work, and directions for further research, including
hybrid architectural models, automatic resource scaling, and the impact of
serverless technologies on the operational costs of cloud applications.

\vspace{0.5cm}

\noindent \textbf{Keywords:} Computer Science, Microservices, Monolith, Cost Analysis, GCP, gRPC, Web Application

\vspace{0.5cm}

\noindent \textbf{OECD Field of Science and Technology Classification:} Computer and Information Sciences